% \section{Kramers--Moyal expansion}\label{app:KM-expansion}
% \noindent
% Just as the integral form~\eqref{eq:n} of $n(\x,t)$ can 
%  be expressed in
% differential form as the continuity equation
% $\partial_t\sp n+\nabla\cdot(n\sp\u)=0$, the integral
% form~\eqref{eq:xi-homo} for $\xi_n(\r,t)$
% can be expressed as a differential equation, 
% taking the form of a Kramers--Moyal expansion.
% There are several ways this equation can be obtained and
% here we derive it following the characteristic-function approach
% of \cite{Risken1989FokkerPlanck}.

% Consider an expansion of Eq.~\eqref{eq:xi-homo} for a small time interval
% (which we take to zero at the end of the calculation).
% To do this it helps to write Eq.~\eqref{eq:xi-homo} using
% arbitrary times $t$ and $t'$ so
% \be\label{eq:xi-homo-alt}
% \xi_n(\r',t')=\int_{\r}p(\r',t'\Mid\r,t)\, \xi_n(\r,t)\,,
% \ee
% Now we will expand $p$ in moments
% \be
% M_{i_1\ldots i_m}(\r;t',t)
% \equiv \int_{\r'} (\r'-\r)_{i_1}\cdots(\r'-\r)_{i_m}\ssp p(\r',t'\Mid\r,t).
% \ee
% (This is closely related to the LPT displacement moments,
% e.g.\ for the first moment ${M}_i=\llangle\Delta\psi_i\rrangle$.)
% % The transition probability can be expanded in terms of
% % the moments $M_{i_1\ldots i_n}(\r';t,t')
% % =\llangle\Delta_{i_1}\cdots\Delta_{i_n}\rrangle$,
% %  $\Delta_i=r_i-r_i'$:
% % \bea
% % M_{i_1\ldots i_n}(\r';t,t')
% % % &\equiv\llangle\Delta_{i_1}\cdots\Delta_{i_n}\rrangle \\
% % &=\int_\r (\r-\r')_{i_1}\cdots(\r-\r')_{i_n}\ssp p(\r,t\Mid\r',t').
% % \eea
% % When inserted back into Eq.~\eqref{eq:xi-homo} we have the
% % Kramers--Moyal expansion. This is useful for obtaining
% % the differential equation form of Eq.~\eqref{eq:xi-homo}, i.e.\
% % in the limit of infinitesimally small time steps.
% % The following calculation can be found in
% % \cite{Risken1989FokkerPlanck}.
% These moments are encoded in the
% generating function~\eqref{eq:pt-r} as
% $Z(\k;\r)=\llangle e^{\im\k\cdot\r_{t'}}\rrangle
% =e^{\im\k\cdot\r}\llangle e^{\im\k\cdot\bm(\r_{t'}-\r)}\rrangle$.
% Expanding the second exponential in a power series,
% the transition probability~\eqref{eq:pt-r} becomes
% \be\nonumber
% p(\r',t'\Mid\r,t)
% =\int_\k e^{-\im\k\cdot(\r'-\r)}
%     \left(1+\sum_{m=1}^\infty\frac{\im^m}{m!} k_{i_1}\cdots k_{i_m}
%         M_{i_1\ldots i_m}(\r;t',t)
%     \right)
% \ee
% Evaluating the integral over $\k$ using the identity
% \be\nonumber
% \int_\k e^{-\im\k\cdot(\r'-\r)}\sp\im^m\sp k_{i_1}\cdots k_{i_m}
% =(-1)^n\ssp\frac{\partial}{\partial r_{i_1}'}
%     \cdots\frac{\partial}{\partial r_{i_m}'}\,\delD(\r'-\r)
% \ee
% and using $\delD(\r'-\r)\sp M(\r)=\delD(\r'-\r)\sp M(\r')$, we have
% \bea
% &p(\r',t'\Mid\r,t) \\
% &=\left(1+\sum_{m=1}^\infty\frac{(-1)^m}{m!}
%     \frac{\partial}{\partial r_{i_1}'}
%     \cdots\frac{\partial}{\partial r_{i_m}'} M_{i_1\ldots i_m}(\r';t',t)
%     \right)\delD(\r'-\r). \nonumber
% \eea
% Inserted back into Eq.~\eqref{eq:xi-homo-alt} and bringing the
% zeroth-order term over to the left-hand side, we have
% \bea\nonumber
% &\xi_n(\r',t')-\xi_n(\r',t) 
% =\sum_{m=1}^\infty\frac{(-1)^m}{m!}
%     \frac{\partial}{\partial r_{i_1}'}
%     \cdots\frac{\partial}{\partial r_{i_m}'} M_{i_1\ldots i_m}(\r';t',t)\ssp
%     \xi_n(\r',t).\nonumber
% \eea
% Expanding both sides, assuming a small time interval $\Delta t=t'-t$,
% the left-hand side reduces to
% $\partial_t\ssp\xi_n\ssp\Delta t+O(\Delta t^2)$;
% on the right-hand side we have
% \be\nonumber
% M_{i_1\ldots i_m}(\r';t+\Delta t,t)
% =K_{i_1\ldots i_m}(\r',t)\cdot\Delta t+O(\Delta t^2)\,,
% \ee
% where
% \be
% K_{i_1\ldots i_m}(\r,t)
% \equiv\lim_{\Delta t\to 0} M_{i_1\ldots i_m}(\r;t+\Delta t,t)/\Delta t
% \ee
% is the Kramers--Moyal coefficient~\citep{Risken1989FokkerPlanck}.
% Inserting these expansions into the foregoing equation,
% dividing both sides by $\Delta t$ and then taking the limit
% $\Delta t\to0$, we obtain
% \be\label{eq:KM-expansion}
% \partial_t\ssp\xi_n(\r,t)
% =\sum_{m=1}^\infty\frac{(-1)^m}{m!}
%     \frac{\partial}{\partial r_{i_1}}
%     \cdots\frac{\partial}{\partial r_{i_m}}
%         \big[K_{i_1\ldots i_m}(\r,t)\, \xi_n(\r,t)\big]\,.
% \ee
% (For simplicity we have dropped the prime on the
% arbitrary separation $\r'$.)
% This is the Kramers--Moyal expansion.

% The pair conservation equation is obtained by 
% neglecting all but the first term  on the right-hand side
% of Eq.~\eqref{eq:KM-expansion}:
% \be
% \partial_t\ssp \xi_n
% +\bm\nabla\cdot(\xi_n\ssp \u_{12})=0\,,
% \ee
% where $\u_{12}(\r,t)=\mbf{K}(\r,t)$ is the 
% pairwise relative velocity~\citep{Davis:1977,Peebles_LSS}.

% % A Fokker--Planck equation is obtained by 
% % neglecting all but the first two terms on the right-hand side 
% % of Eq.~\eqref{eq:KM-expansion}:
% % \be
% % \partial_t\ssp \xi_n
% % =-\frac{\partial}{\partial r_i}({D}_i\ssp\xi_n)
% % +\frac12\frac{\partial}{\partial r_i}
% %     \frac{\partial}{\partial r_j}(D_{ij}\ssp\xi_n)\,.
% % \ee




\section{Derivation of the pair conservation equation}\label{app:KM-expansion}
\noindent
To derive the pair conservation equation~\eqref{eq:pairwise-eqn} from the
the pushforward~\eqref{eq:xi-homo},
%Just as the integral form~\eqref{eq:n} of $n(\x,t)$ can 
%be expressed in differential form as the continuity equation,
%the integral form~\eqref{eq:xi-homo} for $\xi_n(\r,t)$
%can be expressed as a differential equation, 
%known as the pair conservation equation.
%This is derived as follows.
% The derivation we fobelow is essentially the same as the derivation of the
% Kramers--Moyal expansion using the characteristic function
% \cite{Risken1989FokkerPlanck}.
consider an expansion of Eq.~\eqref{eq:xi-homo} over a small time interval
$\Delta t$. We take $\Delta t$ to zero at the end of the calculation.
First it will be convenient to call the initial time 
$t_0$ and write Eq.~\eqref{eq:xi-homo} as
\be\label{eq:xi-homo-alt}
\xi_\rho(\r,t)=\int\dif^3{\q}\; p(\r,t\Mid\q,t_0)\, \xi_\rho(\q,t_0)\,,
\ee
where
$p(\r,t\Mid\q,t_0)=\llangle\delD(\r-\r_{t})\rrangle$.
(Note that that the pair trajectory $\r_{t}$ is a random variable
that depends on $\q_1$ and $\q_2$, and only on
$\q=\q_1-\q_2$ after averaging.)
Now it is convenient to use the Fourier representation of
the delta function to write
\be\nonumber
p(\r,t\Mid\q,t_0)
=\int\frac{\dif^3\k}{(2\pi)^3}\, e^{-\im\k\cdot\r}
    \llangle e^{\im\k\cdot\r_{t}}\rrangle
=\int\frac{\dif^3\k}{(2\pi)^3}\, e^{-\im\k\cdot(\r-\q)}
    \llangle e^{\im\k\cdot\Delta\bm\psi}\rrangle\,,
\ee
where we recalled the relative displacement
$\Delta\bm\psi=\r_t-\q$.
Expand the second exponential in a power series:
\be\label{eq:pt-expansion}
p(\r,t\Mid\q,t_0)
=\int\frac{\dif^3\k}{(2\pi)^3}\, e^{-\im\k\cdot(\r-\q)}
    \Big(1+\im \k\cdot\mbf{D}(\q;t,t_0)+\cdots
    \Big)\,.
\ee
As in the main text, here we defined
$\mbf{D}(\q;t,t_0)
% \equiv\langle\!\langle \r_{t}-\q\rangle\!\rangle
\equiv\langle\!\langle\Delta\bm\psi\rangle\!\rangle$, i.e.\ the expected
relative displacement between time $t_0$ and $t>t_0$.
This is related to the expected separation
$\R_t(\q)\equiv\langle\!\langle\r_t\rangle\!\rangle
=\q+\mbf{D}(\q;t,t_0)$.
(For our particular
transition probability the higher-order
terms will turn out be irrelevant when we take the limit $t\to t_0$,
hence our showing only the leading term.%
\footnote{More specifically,
the higher-order terms,
$\langle\!\langle\Delta\psi_i\sp\Delta\psi_j\rangle\!\rangle$, 
$\langle\!\langle\Delta\psi_i\sp\Delta\psi_j\sp\Delta\psi_k\rangle\!\rangle$, etc,
% \bea
% \langle\!\langle\Delta\psi_i\sp\Delta\psi_j\rangle\!\rangle\,,
% \quad
% \langle\!\langle\Delta\psi_i\sp\Delta\psi_j\sp\Delta\psi_k\rangle\!\rangle\,, \quad\text{etc.} 
% \eea
may be neglected because the relative displacement
$\Delta\bm\psi=(\v_1-\v_2)\sp\Delta t+O(\Delta t^2)$ is at
lowest order $\Delta t$, meaning that the second moment is
$O(\Delta t^2)$, the third moment is $O(\Delta t^3)$, and so on. 
Upon dividing through by $\Delta t$ and taking the
limit $\Delta t\to0$,
these moments vanish, leaving only 
$D_i=\langle\!\langle\Delta\psi_i\rangle\!\rangle$.
})
Now using the identity
$(2\pi)^{-3}\int\dif^3\k\, e^{-\im\k\cdot(\r-\q)}\im\k
=-\nabla_\r\,\delD(\r-\q)$, we have
\be\label{eq:pt-LO}
p(\r,t\Mid\q,t_0)
=\delD(\r-\q)
    -\nabla_\r\cdot
        \big[\mbf{D}(\r;t,t_0)\,\delD(\r-\q)\big]+\cdots
\ee
Here, using the property
$\delD(\r-\q)\sp f(\q)=\delD(\r-\q)\sp f(\r)$, we evaluated
$\mbf{D}$ at separation $\r$.
Inserted back into Eq.~\eqref{eq:xi-homo-alt} and bringing the
zeroth-order term over to the left-hand side, we have
\bea\nonumber
\xi_n(\r,t)-\xi_n(\r,t_0) 
=-\nabla_\r\cdot
        \big[\mbf{D}(\r;t,t_0)\,\xi_n(\r,t_0)\big]+\cdots
\eea
Now write $t=t_0+\Delta t$ and expand both sides about the
initial time $t_0$, assuming a small time interval $\Delta t$.
On the left-hand side,
$\xi_n(\r,t)=\xi_n(\r,t_0)+\partial_{t}\ssp\xi_n(\r,t)|_{t=t_0}\Delta t
+O(\Delta t^2)$.
On the right-hand side,
\be\label{eq:M}
\mbf{D}(\r;t_0+\Delta t,t_0)
\equiv\llangle\Delta\bm\psi\rrangle
=\frac1a\v_{12}(\r,t_0)\sp\Delta t+O(\Delta t^2)
\ee
where we have inserted the relative displacement
$\Delta\bm\psi=\r_t-\q=\frac1a(\v_1-\v_2)\sp\Delta t+O(\Delta t^2)$
with $\v_1=\v(\q_1)$ and $\v_2=\v(\q_2)$,
and recognized the pairwise streaming velocity, defined as
$\v_{12}=\llangle\v_1-\v_2\rrangle$.
Here the scale factor enters because $\v_1=a\sp\dif\x_1/\dif t$ and
$\v_2=a\sp\dif\x_2/\dif t$ are the physical velocities
conventionally used.
The expression now reads
\be\label{eq:pair-exp}
\partial_t\ssp\xi_n(\r,t)\,\Delta t
=-\frac1a\nabla_\r\cdot
        \big[\xi_n(\r,t)\,\v_{12}(\r;t)\big]\sp\Delta t
    +O(\Delta t^2)\,.
\ee
(Since the final time $t$ no longer appears in this local-in-time
expression, we have relabelled the arbitrary initial time $t_0\to t$.)
Upon dividing both sides by $\Delta t$ and taking the limit
$\Delta t\to0$,
we recover the pair conservation equation~\eqref{eq:pairwise-eqn}.
%\be\label{eq:pair-eqn}
%\partial_t\ssp\xi_n
%+\frac1a\bm\nabla_\r\cdot
%    \big(\xi_n\,\v_{12}\big)=0\,.
%\ee
%This is the pair conservation equation~\citep{Davis:1977}.
%There are also equations for the conservation of triplets, quadruplets, etc, all of which have the same basic form as
%Eq.~\eqref{eq:pair-eqn}
%and can be derived following a similar calculation to that above.
%For explicit formula and their derivation from the BBGKY hierarchy,
%see \cite{Peebles_LSS}, section 71.B.


%Finally note that Eq.~\eqref{eq:pair-eqn} is generally not closed.
%Since to solve for the zeroth
%moment $\xi_n$ one needs to solve for the first moment
%$\u_{12}$, which in turn requires solving for the second moment
%(the pairwise velocity dispersion), and so on.
%This requires partial solution of the full BBGKY hierarchy.
%Thus a solution of $\xi_n$ generally depends not just on
%$\u_{12}$ but also higher moments.








% \section{Derivation of the pair conservation equation from the transition probability}\label{app:KM-expansion}
% \noindent
% Just as the integral form~\eqref{eq:n} of $n(\x,t)$ can 
%  be expressed in
% differential form as the continuity equation
% $\partial_t\sp n+\nabla\cdot(n\sp\u)=0$, the integral
% form~\eqref{eq:xi-homo} for $\xi_n(\r,t)$
% can be expressed as a differential equation, 
% taking the form of a Kramers--Moyal expansion.
% There are several ways this equation can be obtained and
% here we derive it following the characteristic-function approach
% of \cite{Risken1989FokkerPlanck}.

% Consider an expansion of Eq.~\eqref{eq:xi-homo} for a small time interval
% (which we take to zero at the end of the calculation).
% To do this it will be convenient to call the initial time 
% $t$ and the final time $t'$. We also
% relabel $\q\to\r$ and $\r\to\r'$. Equation~\eqref{eq:xi-homo} 
% then reads
% \be\label{eq:xi-homo-alt}
% \xi_n(\r',t')=\int_{\r}p(\r',t'\Mid\r,t)\, \xi_n(\r,t)\,,
% \ee
% where
% $p(\r',t'\Mid\r,t)=\langle\!\langle\delD(\r'-\r_{t'})\rangle\!\rangle$.
% (Note that $\r_{t'}$ is a stochastic variable that depends on
% $\r$ after averaging.)
% Now expand the delta function in plane waves as
% \be\nonumber
% p(\r',t'\Mid\r,t)
% =\int_\k e^{-\im\k\cdot\r'}
%     \langle\!\langle e^{\im\k\cdot\r_{t'}}\rangle\!\rangle
% =\int_\k e^{-\im\k\cdot(\r'-\r)}
%     \langle\!\langle e^{\im\k\cdot(\r_{t'}-\r)}\rangle\!\rangle
% \ee
% Expanding the second exponential in a power series,
% \be\nonumber
% p(\r',t'\Mid\r,t)
% =\int_\k e^{-\im\k\cdot(\r'-\r)}
%     \Big(1+\im \k\cdot\mbf{M}(\r;t',t)+\cdots
%     \Big)\,.
% \ee
% Here we defined $\mbf{M}(\r;t',t)
% =\langle\!\langle \r_{t'}-\r\rangle\!\rangle$. For our particular
% transition probability the higher-order
% terms will turn out be irrelevant when we take the limit $t\to t'$.
% Hence we exhibit only the leading-order term.
% Now using the identity
% $\int_\k e^{-\im\k\cdot(\r'-\r)}\im\k
% =-\frac{\partial}{\partial\r'}\sp\delD(\r'-\r)$, we have
% \be\nonumber
% p(\r',t'\Mid\r,t)
% =\delD(\r'-\r)
%     -\frac{\partial}{\partial\r'}\cdot
%         \big[\mbf{M}(\r';t',t)\,\delD(\r'-\r)\big]+\cdots
% \ee
% Here we have also used the property
% $\delD(x-y)\sp f(y)=\delD(x-y)\sp f(x)$.
% Inserted back into Eq.~\eqref{eq:xi-homo-alt} and bringing the
% zeroth-order term over to the left-hand side, we have
% \bea\nonumber
% \xi_n(\r',t')-\xi_n(\r',t) 
% =-\frac{\partial}{\partial\r'}\cdot
%         \big[\mbf{M}(\r';t',t)\,\xi_n(\r',t)\big]+\cdots
% \eea
% Now write $t'=t+\Delta t$ and expand both sides about the
% initial time $t$, assuming a small time interval $\Delta t$.
% We have
% $\xi_n(\r',t')=\xi_n(\r',t)+\partial_t\sp\xi_n(\r',t)\sp\Delta t
% +O(\Delta t^2)$.
% We also have
% $\mbf{M}(\r';t',t)
% =\u_{12}(\r',t)\sp\Delta t+O(\Delta t^2)$,
% where $\u_{12}=\langle\!\langle\u_1-\u_2\rangle\!\rangle$ is the
% pairwise streaming velocity.
% Our expression now reads
% \be
% \partial_t\ssp\xi_n(\r',t)\,\Delta t
% =-\frac{\partial}{\partial\r'}\cdot
%         \big[\xi_n(\r',t)\,\u_{12}(\r';t)\big]\sp\Delta t
%     +O(\Delta t^2)
% \ee
% Dividing both sides by $\Delta t$, taking the limit $\Delta t\to0$,
% and relabelling $\r'\to\r$, we get
% \be
% \partial_t\ssp\xi_n
% +\bm\nabla_\r\cdot
%     \big(\xi_n\,\u_{12}\big)=0\,.
% \ee
% This is the pair conservation equation~\citep{Davis:1977,Peebles_LSS}.

% We mentioned that the higher-order moments of the transition
% probability, e.g.\
% \bea
% M_{ij}(\r;t',t)
% &\equiv\int_{\r'} (\r'-\r)_i (\r'-\r)_j \sp p(\r',t'\Mid\r,t) \\
% &=\langle\!\langle(\r_{t'}-\r)_i(\r_{t'}-\r)_j\rangle\!\rangle
% \eea
% may be neglected. This is because when we insert the perturbative
% solution $\r_{t'}=\r+(\u_1-\u_2)\sp\Delta t+O(\Delta t^2)$
% it is not hard to see that $M_{ij}=O(\Delta t^2)$. Hence
% when we divide by $\Delta t$ and take the limit $\Delta t\to0$
% this moment vanishes.
% Indeed all moments higher than the first vanish in this limit,
% e.g.\ for third moment $M_{ijk}=O(\Delta t^3)$, and so on.






% Here's a more direct derivation. Write
% \be\nonumber
% p_t(\x_1,\x_2\Mid\q_1,\q_2)
% =\langle\!\langle\delD[\x_1-\x_t(\q_1)]\ssp\delD[\x_2-\x_t(\q_2)] 
%     \rangle\!\rangle\,,
% \ee
% Consider the transition of a small time interval, $t=\Delta t$.
% Expand each trajectory about the initial time:
% \bea
% \x_t(\q_1)
% &=\q_1+\frac{\dif\x_t(\q_1)}{\dif t}\bigg|_{t=0}\,\Delta t
%     +O(\Delta t^2)\,,\nonumber\\
% \x_t(\q_2)
% &=\q_2+\frac{\dif\x_t(\q_2)}{\dif t}\bigg|_{t=0}\,\Delta t
%     +O(\Delta t^2)\,,
% \eea
% Now we expand the transition probability:
% \bea
% &p_t(\x_1,\x_2\Mid\q_1,\q_2) 
% =\delD(\x_1-\q_1)\ssp\delD(\x_2-\q_2) \nonumber\\
% &
%     -\left\langle\!\left\langle\frac{\dif\x_t(\q_1)}{\dif t}\bigg|_{t=0}\cdot\frac{\partial}{\partial\x_1}\ssp\delD(\x_1-\q_1)\ssp\delD(\x_2-\q_2)+(1\leftrightarrow2)\right\rangle\!\right\rangle\Delta t \nonumber\\
% &
%  + O(\Delta t^2)
%     \nonumber
% \eea
% This yields
% \bea
% \partial_t\ssp\xi_n(\x_1,\x_2,t)\ssp\Delta t
% &=-\bm\nabla_{\x_1}\cdot\big[
%     \langle\!\langle\u_1\rangle\!\rangle\ssp
%         \xi_{n0}(\x_1,\x_2)\big]\Delta t \nonumber\\
% &\quad-\bm\nabla_{\x_2}\cdot\big[
%     \langle\!\langle\u_2\rangle\!\rangle\ssp
%         \xi_{n0}(\x_1,\x_2)\big]\Delta t
% +O(\Delta t^2)\nonumber
% \eea
% Using that the correlations depend on $\x_1-\x_2$ (homogeneity)
% we obtain, in the limit $\Delta t\to0$,
% \be
% \partial_t\ssp\xi_n+\bm\nabla_\r\cdot\big(\xi_n\ssp\u_{12}\big)=0
% \ee



